% Point to the main file
\documentclass[../Main.tex]{subfiles}

% Start the document
\begin{document}

\begin{alphabet}
    \item 
    {
        We can split this into 

        \begin{center}
            total = 0 numerals + 1 numeral + at least 2 numerals
        \end{center}

        \begin{center}
            at least 2 numerals = total - 0 numerals - 1 numeral
        \end{center}

        Let's start by finding the total. Since there are no restrictions, each position has 36 options

        \begin{align*}
            total &= \underline{36} \cdot \underline{36} \cdot \underline{36} \cdot \underline{36} \cdot \underline{36} \cdot \underline{36} \cdot \underline{36} = 36^7 
        \end{align*}

        Now for 0 numerals, the restriction is that each character has to be an uppercase letter, so each position has 26 options

        \begin{align*}
            0 \text{ numerals} &= \underline{26} \cdot \underline{26} \cdot \underline{26} \cdot \underline{26} \cdot \underline{26} \cdot \underline{26} \cdot \underline{26} = 26^7
        \end{align*}

        Lastly, 1 exact numeral means only one slot has 10 options, but out of 7 slots. So we multiply the count for 1 numeral by the ways we can put that one numeral. The rest of the slots MUST be letters

        \begin{align*}
            1 \text{ numeral} = {7 \choose 1} \cdot \underline{10} \cdot \underline{26} \cdot \underline{26} \cdot \underline{26} \cdot \underline{26} \cdot \underline{26} \cdot \underline{26} &= 7 \cdot 10 \cdot 26^6 \\
            &= 70 \cdot 26^6
        \end{align*}

        Using our equation from above,

        \begin{align*}
            \text{at least 2 numerals} = \boxed{6^7 - 26^7 - 70 \cdot 26^6}
        \end{align*}
    }
    
    \item
    {
        First, let's find out the number of license plates with MILES being at the start

        \begin{align*}
            \underbrace{1}_M \cdot \underbrace{1}_I \cdot \underbrace{1}_L \cdot \underbrace{1}_E \cdot \underbrace{1}_S \cdot 36 \cdot 36 = 36^2
        \end{align*}

        Now let's see how many times we can arrange MILES with 7 slots

        \begin{enumerate}
            \item MILESxx
            \item xMILESx
            \item xxMILES
        \end{enumerate}

        Since there's three ways to putting MILES into 7 characters, we can just multiple $36^2$ by 3 to get

        $$\boxed{3 \cdot 36^2}$$
    }

    \item
    {
        We can think of this as:

        \begin{enumerate}
            \item choose 3 different characters out of 36
            \item choose 3 different spots of 7 for the first character
            \item choose 3 spots out of 4 for the second character
            \item choose 1 spot out of 1 for the last character
        \end{enumerate}

        Mathematically, this ends up being

        \begin{align*}
            \boxed{{36 \choose 3} \cdot {7 \choose 3} \cdot {4 \choose 3} \cdot {1 \choose 1}}
        \end{align*}
    }

    \item 
    {
        Like the last one, we can think of this as

        \begin{enumerate}
            \item choose 3 different numerals out of 10
            \item choose 1 letter out of 26
            \item choose 3 spots out of 7 to place the numerals 
            \item permute the 3 numbers
            \item choose 4 spots of out 4 to place the lettters (basically just put the letters where the numbers aren't)
        \end{enumerate}

        This ends up being

        \begin{align*}
            & {10 \choose 3} \cdot {26 \choose 1} \cdot {7 \choose 3} \cdot \prescript{}{3}{P}_3 \cdot {4 \choose 4} \\
            &= \boxed{{10 \choose 3} \cdot {26 \choose 1} \cdot {7 \choose 3} \cdot 3!} \\
        \end{align*}
    }

    \item
    {
        For this one, since alphabetical order means A-Z (1 - 26) and there's 7 letters, the start can only be up to
        
        $$26 - 7 + 1 = 20 \; (\text{T})$$

        We can think of the first character essentially controlling the entire string, so in reality all we need to do is choose
        a valid first letter.

        Using the range

        $$\left[1, 20 \right]$$

        We have 20 letters to choose from, means our answer is just

        $$\boxed{20}$$
    }

\end{alphabet}

% End the document
\end{document}